\section*{Hoofdstuk \ref{ch:g11:infrared} --- Infraroodspectroscopie}

\begin{solution}{exo:g11:infrared:1}
$\lambda = 1/3300 = \qty{3.03e-4}{cm} = \qty{3.03}{\micro m}$;
$\lambda = 1/1050 = \qty{9.52e-4}{cm} = \qty{9.52}{\micro m}$. De band bij
\qty{3300}{cm^{-1}}, met het grotere \omterm{def:g11:infrared:wavenumber}{golfgetal}, hoort bij de snellere
trilling.
\end{solution}

\begin{solution}{exo:g11:infrared:2}
$\qty{10}{\micro m} = \qty{1.0e-3}{cm}$, dus $\sigma = 1/\num{1.0e-3} =
\qty{1000}{cm^{-1}}$.
\end{solution}

\begin{solution}{exo:g11:infrared:3}
De \ce{C=O} van een \omterm{def:g11:functional-groups:carbonyl}{keton}; een alcohol-\ce{O-H} met \omterm{def:g11:polarity-and-cohesion:hydrogen-bond}{waterstofbruggen};
\ce{C-H}-bindingen van een koolstofketen.
\end{solution}

\begin{solution}{exo:g11:infrared:4}
De band bij \qty{1715}{cm^{-1}}: de \ce{C=O}-binding.
\end{solution}

\begin{solution}{exo:g11:infrared:5}
Het instrument meet het licht dat door het monster gaat; waar een binding
absorbeert, gaat er minder licht door en zakt de kromme. $T =
\qty{100}{\percent}$ betekent dat er niets van het licht met dat \omterm{def:g11:infrared:wavenumber}{golfgetal}
wordt opgenomen.
\end{solution}

\begin{solution}{exo:g11:infrared:6}
De eerste heeft een \ce{C=O}-band bij \qty{1715}{cm^{-1}} en geen
\ce{O-H}-band. Het is een \omterm{def:g11:functional-groups:carbonyl}{keton}, butanon (ook
butaan-2-on genoemd), met de formule \ce{CH3-CO-CH2-CH3}. De tweede heeft een \ce{O-H}-band en geen
\ce{C=O}: een \omterm{def:g11:functional-groups:alcohol}{alcohol}, bijvoorbeeld butaan-1-ol.
\end{solution}

\begin{solution}{exo:g11:infrared:7}
De \ce{O-H}-groepen van het zuur zijn heel sterk door \omterm{def:g11:polarity-and-cohesion:hydrogen-bond}{waterstofbruggen}
gebonden (de \omterm{def:g7:atoms-and-molecules:molecule}{moleculen} vormen paren, elke \ce{O-H} gebonden aan de
\ce{C=O} van de andere), en elk \omterm{def:g7:atoms-and-molecules:molecule}{molecuul} is een beetje anders gebonden:
de band wordt nog losser en spreidt zich over een heel breed gebied.
\end{solution}

\begin{solution}{exo:g11:infrared:8}
In de vloeistof (A) is de \ce{O-H}-band breed, bij
\qty{3350}{cm^{-1}}; in het gas (B), waar de \omterm{def:g7:atoms-and-molecules:molecule}{moleculen} ver uit elkaar
zitten en geen \omterm{def:g11:polarity-and-cohesion:hydrogen-bond}{waterstofbruggen} vormen, is het een smalle band bij
\qty{3666}{cm^{-1}}.
\end{solution}

\begin{solution}{exo:g11:infrared:9}
Nauwelijks: bij \qty{1730}{cm^{-1}} voor het \omterm{def:g11:functional-groups:carbonyl}{aldehyde} en
\qty{1715}{cm^{-1}} voor het \omterm{def:g11:functional-groups:carbonyl}{keton}, waarden die dicht bij elkaar liggen.
Een vergelijking van het hele spectrum, met het \omterm{def:g11:infrared:band}{vingerafdrukgebied}, met
referentiespectra zou beslissen, of een andere techniek.
\end{solution}

\begin{solution}{exo:g11:infrared:10}
Een \omterm{def:g11:functional-groups:carboxylic-acid}{ester} (\ce{C=O} bij \qty{1735}{cm^{-1}}, sterke \ce{C-O}, geen
\ce{O-H}), bijvoorbeeld ethylethanoaat. Ethaanzuur is \ce{C2H4O2}, niet
\ce{C4H8O2}. Butaanzuur heeft de juiste formule maar zou de heel brede
\ce{O-H}-band van zuren laten zien: uitgesloten.
\end{solution}

\begin{solution}{exo:g11:infrared:11}
Propaan-1-ol is een \omterm{def:g11:functional-groups:alcohol}{alcohol}: zoals A. Propaanzuur is een \omterm{def:g11:functional-groups:carboxylic-acid}{carbonzuur}:
zoals D.
\end{solution}

\begin{solution}{exo:g11:infrared:12}
De brede \ce{O-H}-band bij \qty{3350}{cm^{-1}} krimpt, terwijl een sterke
\ce{C=O}-band bij \qty{1715}{cm^{-1}} groeit. De reactie is voltooid als
de \ce{O-H}-band verdwenen is en de \ce{C=O}-band niet meer groeit.
\end{solution}

\begin{solution}{exo:g11:infrared:13}
Butaanzuur vertoont een heel brede \ce{O-H}-band van 2500 tot
\qty{3300}{cm^{-1}}, de \omterm{def:g11:functional-groups:carboxylic-acid}{ester} geen; de \ce{C=O} van het zuur ligt bij
\qty{1710}{cm^{-1}}, breder, die van de \omterm{def:g11:functional-groups:carboxylic-acid}{ester} bij \qty{1735}{cm^{-1}},
smal, met een sterke \ce{C-O}-band bij \qty{1240}{cm^{-1}}. Het zuur
heeft een \ce{O-H}-groep met \omterm{def:g11:polarity-and-cohesion:hydrogen-bond}{waterstofbruggen}, die ook zijn \ce{C=O} een
beetje losser maakt.
\end{solution}

\begin{solution}{exo:g11:infrared:14}
Overgebleven ethanol (de \omterm{def:g11:functional-groups:alcohol}{alcohol} waaruit de \omterm{def:g11:functional-groups:carboxylic-acid}{ester} gemaakt wordt) of water
(uit de \omterm{def:g4:air-a-mixture-of-gases:air}{lucht} of van het wassen): beide hebben \ce{O-H}-groepen met
\omterm{def:g11:polarity-and-cohesion:hydrogen-bond}{waterstofbruggen}. Je kunt het monster vergelijken met een
referentiespectrum, drogen, destilleren, of het kookpunt meten.
\end{solution}

\begin{solution}{exo:g11:infrared:15}
Een \omterm{def:g10:lewis-and-shape:multiple-bond}{dubbele binding} is een stijvere veer dan een enkele binding tussen
dezelfde \omterm{def:g7:atoms-and-molecules:atom}{atomen}: ze trilt sneller, bij een hoger \omterm{def:g11:infrared:wavenumber}{golfgetal} (1715 tegen
\qty{1050}{cm^{-1}}). De bindingen \ce{O-H}, \ce{N-H} en \ce{C-H} bevatten
allemaal een waterstofatoom, het lichtste van alle \omterm{def:g7:atoms-and-molecules:atom}{atomen}: lichte
balletjes aan een veer trillen snel, vandaar \omterm{def:g11:infrared:wavenumber}{golfgetallen} rond
2900--\qty{3600}{cm^{-1}}.
\end{solution}

\begin{solution}{pb:g11:infrared:1}
\textbf{1.} \ce{CH3-CH2-OH}; \ce{CH3-CO-CH3}; \ce{CH3-COOH}.

\textbf{2.} \ce{C2H6O}; \ce{C3H6O}; \ce{C2H4O2}.

\textbf{3.} Hydroxyl; carbonyl (\omterm{def:g11:functional-groups:carbonyl}{keton}); carboxyl.

\textbf{4.} Ethanol: \ce{O-H} en \ce{C-H}. Propanon: \ce{C-H} en
\ce{C=O}. Ethaanzuur: \ce{O-H}, \ce{C-H} en \ce{C=O}.

\textbf{5.} Geen \ce{C=O}-band, een brede \ce{O-H}-band: ethanol.

\textbf{6.} Een heel brede \ce{O-H}-band van 2500 tot \qty{3300}{cm^{-1}}
en een sterke \ce{C=O} bij \qty{1710}{cm^{-1}}: ethaanzuur.

\textbf{7.} Propanon: de heel sterke \ce{C=O}-band bij
\qty{1715}{cm^{-1}}, zonder \ce{O-H}.

\textbf{8.} Waarschijnlijk wel (azijn, \omterm{def:g6:solutions-and-solubility:solute}{oplosmiddel}, \omterm{def:g11:functional-groups:alcohol}{alcohol}), maar aan
onbekende vloeistoffen ruiken is gevaarlijk; het spectrum is veilig en
zeker.

\textbf{9.} De \ce{O-H}-groepen van het zuur zijn sterker door
\omterm{def:g11:polarity-and-cohesion:hydrogen-bond}{waterstofbruggen} gebonden (paren van \omterm{def:g7:atoms-and-molecules:molecule}{moleculen}), dus de band ligt lager
en is veel breder.

\textbf{10.} \ce{CH3-O-CH3}, \ce{C2H6O}; het heeft geen \ce{O-H}, dus geen
\ce{O-H}-band.

\textbf{11.} De \ce{C=O}-band: bij \qty{1730}{cm^{-1}} voor propanal,
bij \qty{1715}{cm^{-1}} voor propanon.

\textbf{12.} De \ce{O-H}-band: een \omterm{def:g11:functional-groups:carboxylic-acid}{ester} heeft geen \ce{O-H}. Zijn
\ce{C=O} zou bij \qty{1735}{cm^{-1}} liggen.

\textbf{13.} \omterm{def:g11:organic-skeletons:isomer}{Isomeren} hebben dezelfde formule; hun \omterm{def:g11:functional-groups:functional-group}{karakteristieke groepen}
verschillen, en het spectrum laat de groepen zien.

\textbf{14.} \qty{1715}{cm^{-1}}.

\textbf{15.} $1715 \times 100 = \qty{1.715e5}{m^{-1}}$.

\textbf{16.} $\lambda = 1/\num{1.715e5} = \qty{5.83e-6}{m} =
\qty{5.83}{\micro m}$.

\textbf{17.} Nee: zichtbaar licht loopt van ongeveer 0.4 tot
\qty{0.75}{\micro m}. Het is infraroodstraling.

\textbf{18.} $\num{1.05e5}\ \si{m^{-1}}$, dus $\lambda = \qty{9.52e-6}{m}
= \qty{9.52}{\micro m}$.

\textbf{19.} De \ce{C=O}-binding (hoger \omterm{def:g11:infrared:wavenumber}{golfgetal}, kortere golflengte).

\textbf{20.} Ongeveer \qty{5.8}{\micro m}.
\end{solution}
